\chapter{Gestion des Risques}

\section{Identification et Priorisation}

Conformément au Chapitre 9 du cours (Les Risques), nous avons identifié, priorisé et défini des actions de prévention pour les risques de ce projet. La gestion des risques comprend 4 étapes : \textbf{identifier}, \textbf{prioriser}, \textbf{prévenir} et \textbf{suivre}.

\subsection{Risques Techniques}

\begin{table}[H]
\centering
\begin{tabularx}{\textwidth}{C{0.7cm} X C{1.8cm} C{1.5cm} C{1.8cm}}
\toprule
\textbf{\color{primarycolor}\#} &
\textbf{\color{primarycolor}Risque} &
\textbf{\color{primarycolor}Probabilité} &
\textbf{\color{primarycolor}Impact} &
\textbf{\color{primarycolor}Criticité} \\
\midrule
R1 & Perte de données Kafka (le broker crash) & Moyenne & Élevé & Critique \\
\addlinespace[0.1cm]
R2 & Latence du traitement (topologies trop lentes) & Faible & Moyen & Modéré \\
\addlinespace[0.1cm]
R3 & Incompatibilité des versions Docker/Kafka/Flink & Moyenne & Élevé & Critique \\
\addlinespace[0.1cm]
R4 & Mémoire insuffisante (Docker consomme trop de RAM) & Moyenne & Moyen & Modéré \\
\addlinespace[0.1cm]
R5 & Incohérence des données (Replay $\neq$ Temps réel) & Faible & Élevé & Modéré \\
\bottomrule
\end{tabularx}
\caption{Matrice des risques techniques}
\end{table}

\subsection{Risques de Gestion de Projet}

\begin{table}[H]
\centering
\begin{tabularx}{\textwidth}{C{0.7cm} X C{1.8cm} C{1.5cm} C{1.8cm}}
\toprule
\textbf{\color{primarycolor}\#} &
\textbf{\color{primarycolor}Risque} &
\textbf{\color{primarycolor}Probabilité} &
\textbf{\color{primarycolor}Impact} &
\textbf{\color{primarycolor}Criticité} \\
\midrule
R6 & Retard de livraison (complexité sous-estimée) & Élevée & Élevé & Critique \\
\addlinespace[0.1cm]
R7 & Mauvaise répartition des tâches & Moyenne & Moyen & Modéré \\
\addlinespace[0.1cm]
R8 & Manque de compétences (technologies nouvelles) & Élevée & Moyen & Modéré \\
\addlinespace[0.1cm]
R9 & Communication défaillante MOA/MOE & Faible & Moyen & Faible \\
\bottomrule
\end{tabularx}
\caption{Matrice des risques de gestion de projet}
\end{table}

\section{Plans de Mitigation}

\begin{table}[H]
\centering
\begin{tabularx}{\textwidth}{C{0.7cm} X X}
\toprule
\textbf{\color{primarycolor}\#} &
\textbf{\color{primarycolor}Action Préventive} &
\textbf{\color{primarycolor}Action Corrective} \\
\midrule
R1 & Configuration de la rétention (7 jours) et réplication & Replay depuis le dernier offset connu \\
\addlinespace[0.1cm]
R2 & Partitionnement des topics (3 partitions) & Augmenter les task slots Flink \\
\addlinespace[0.1cm]
R3 & Figer les versions dans docker-compose & Tester avec des versions antérieures \\
\addlinespace[0.1cm]
R6 & Estimation par Planning Poker & Réduire le périmètre (prioriser) \\
\addlinespace[0.1cm]
R8 & Auto-formation et prototypage en Sprint 0 & Pair programming et support mutuel \\
\bottomrule
\end{tabularx}
\caption{Plans de mitigation des risques critiques}
\end{table}

\section{Le Cône d'Incertitude}

Référence au Chapitre 8 du cours (Estimation de charge). Au début du projet, l'incertitude liée à la technologie Kafka était très forte (\textbf{facteur $\times$4}). Au fur et à mesure des Sprints, la maîtrise technique s'est améliorée :

\begin{itemize}
    \item \textbf{Début du projet} : incertitude $\times$4 (technologies inconnues Kafka/Flink)~;
    \item \textbf{Après Sprint 1} : incertitude $\times$2 (l'infrastructure fonctionne)~;
    \item \textbf{Après Sprint 2} : incertitude $\times$1 (les topologies sont validées).
\end{itemize}

\section{Suivi des Risques par Sprint}

\begin{table}[H]
\centering
\begin{tabularx}{\textwidth}{L{3cm} C{3cm} C{3cm} C{3cm}}
\toprule
\textbf{\color{primarycolor}Risque} &
\textbf{\color{primarycolor}Sprint 1} &
\textbf{\color{primarycolor}Sprint 2} &
\textbf{\color{primarycolor}Sprint 3} \\
\midrule
R1 Perte données & Non traité & Rétention configurée & Validé \\
\addlinespace[0.1cm]
R2 Latence & Pas de problème & Acceptable & OK \\
\addlinespace[0.1cm]
R3 Versions & Résolu & -- & -- \\
\addlinespace[0.1cm]
R4 Mémoire & Docker lourd & Optimisé & OK \\
\addlinespace[0.1cm]
R6 Retard & Sprint 1 long & Rattrapage & Livré \\
\bottomrule
\end{tabularx}
\caption{Suivi de l'évolution des risques par sprint}
\end{table}

\begin{warningbox}
    \textbf{\textcolor{red!70!black}{\large Risque Technique Majeur}}\\[0.3cm]
    Les technologies Big Data/Streaming sont complexes à configurer. L'utilisation de \texttt{docker-compose} a permis d'encapsuler cette complexité et d'atténuer le risque de problèmes d'environnement. En figeant les versions (\texttt{cp-kafka:7.3.0}, \texttt{flink:1.17}), nous avons éliminé le risque d'incompatibilité R3.
\end{warningbox}
